\subsection{Joint mass and overtone limit}
\label{sec:lovelock-crossover}

The Airy equation also resolves the regime in which the mass
term in \eqref{eq:lovelock-k-z} remains finite in the inner scaling.
Let \(N=n+h\) as above, take \(\nu\geq1\), and hold
\begin{equation}
 \chi=\frac{\sqrt{\nu^2-1}}{N^{2/3}}
 \label{eq:lovelock-scaled-mass}
\end{equation}
fixed as \(N\to\infty\). We compute the slope at \(k^2=0\).
Since \(\sinh^2 z_s=-1\), the mass contributes a constant
\((\nu^2-1)/(4w^{4/3})\) to the leading inner equation in \(y\).
In the coordinate \(X=e^{\ii\pi/3}y\) of
\eqref{eq:lovelock-perturbed-airy}, the equation becomes
\begin{equation}
 F''-\left(X-\zeta-\frac gX\right) F=0,\qquad
 \zeta=e^{-\ii\pi/3}\frac{\chi^2}{4\,2^{2/3}}.
 \label{eq:lovelock-displaced-airy}
\end{equation}
Here \(w=(1-\ii)N\) at \(k=0\), and \(g\) has the same
definition as in \eqref{eq:lovelock-perturbed-airy}.
The mass displaces the Airy turning point. The momentum
perturbation remains anchored at the singularity \(X=0\).

\subsubsection{The connection coefficient at fixed scaled mass}

Normalize \(F(X;g,\zeta)\) to \(\operatorname{Ai}(X-\zeta)\)
at positive large \(X\), and define
\(a(g,\zeta)=F(0;g,\zeta)\) and
\(a_1(\zeta)=\partial_g a(0,\zeta)\).
Variation of parameters gives
\begin{multline}
 a_1(\zeta)=\pi\int_0^\infty\frac{\operatorname{Ai}(t-\zeta)}t
 \big[\operatorname{Bi}(-\zeta)\operatorname{Ai}(t-\zeta)\\
       -\operatorname{Ai}(-\zeta)\operatorname{Bi}(t-\zeta)\big]\dd t .
 \label{eq:lovelock-shifted-variation}
\end{multline}
The bracket vanishes linearly at the origin, and the product
\(\operatorname{Ai}\operatorname{Bi}\) decays as \(t^{-1/2}\)
at large positive \(t\). With \(\sigma=e^{2\pi\ii/3}\), the
rotated solutions now use
\((X,g,\zeta)\mapsto(\sigma^{-j}X,\sigma^j g,\sigma^{-j}\zeta)\).
Each adjacent Wronskian keeps its \(g=0\) value by the
same large-\(X\) normalization as in section~\ref{sec:lovelock-airy}.
At \(g=0\) the connection relation is the Airy identity
with argument \(X-\zeta\), so \(S_{\rm A}(0,\zeta)=-\sigma^{-1}\).
Taking its \(g\) derivative at the origin gives
\begin{equation}
 D_{\rm A}(\zeta)\equiv\partial_g\log S_{\rm A}(0,\zeta)
 =-\frac{a_1(\zeta)+a_1(\sigma\zeta)+a_1(\sigma^{-1}\zeta)}
         {\operatorname{Ai}(-\zeta)} .
 \label{eq:lovelock-displaced-connection}
\end{equation}
At \(\zeta=0\) this equals \(\Gamma(1/3)^2/3^{1/3}\), recovering
\eqref{eq:lovelock-airy-limit}.

At fixed \(\chi\), the shifted Airy region has size
\(|w|^{-2/3}\) in \(z\). The exterior momentum derivative
is subleading by the same overlap estimate as in the
fixed-mass calculation, since the mass shift in \(z\) is
of this order as well. The matched pole condition
therefore predicts
\begin{equation}
 N^{2/3} b_n\ \longrightarrow\
 \mathcal B(\chi)=
 \frac{\ii e^{-\ii\pi/3}}{8\mathcal Z_{\rm L}(1-\ii)^{2/3}}
 D_{\rm A}(\zeta).
 \label{eq:lovelock-crossover}
\end{equation}
Only the \(k^2\) coefficient of the pole condition is varied
here. Its value at zero momentum follows from the exact spectrum.

\begin{table}[t]
 \centering\small
 \begin{tabular}{cccc}
 \hline
 \(\chi\) & \(\mathcal B(\chi)\) & Error at \(n=64\) & Error at \(n=256\) \\
 \hline
 0 & $0.4293294+0.1150385\ii$ & 1.12e-04 & 1.78e-05 \\
1 & $0.3922517+0.1456450\ii$ & 8.31e-05 & 1.40e-05 \\
2 & $0.2873197+0.1840777\ii$ & 5.25e-05 & 9.50e-06 \\
3 & $0.1811542+0.1651296\ii$ & 5.23e-05 & 1.01e-05 \\
5 & $0.1002772+0.1002742\ii$ & 9.90e-06 & 2.24e-06 \\
\hline

 \end{tabular}
 \caption{Joint mass and overtone limit of the momentum slopes.
 The errors compare \eqref{eq:lovelock-crossover} with
 \(N^{2/3}b_n\) computed from \eqref{eq:lovelock-exact-dispersion}
 at the indicated fixed \(\chi\).}
 \label{tab:lovelock-crossover}
\end{table}

Table~\ref{tab:lovelock-crossover} compares the local connection
calculation with the exact moments of the boundary-value problem.
For each \(n\), the mass parameter in the exact formula is chosen
to satisfy \eqref{eq:lovelock-scaled-mass}. In the Airy integrals
we use the Maclaurin series on a finite interval and integrate
the asymptotic expansion of \(\operatorname{Ai}\operatorname{Bi}\)
in the tail \cite{NIST:DLMF}. Changing the matching radius
from 13 to 16 and increasing the truncation order checks that
this numerical error is below the finite-overtone effect.

\subsubsection{The geometric limit of the momentum derivative}

Take \(\chi\) large after the joint limit. In \(X/\zeta\) the
simple-turning-point problem has large WKB parameter
\(\zeta^{3/2}\), and the distant pole at \(X=0\) stays outside
the connection from the outer oscillatory region to the
displaced turning point. The leading \(g\) derivative of
the WKB action from the turning point to infinity is
\begin{equation}
 -\frac12\int_0^\infty\frac{\dd t}{(\zeta+t)\sqrt t}
    =-\frac{\pi}{2\sqrt\zeta}.
 \label{eq:lovelock-displaced-action}
\end{equation}
The endpoint derivative of the action vanishes at the
turning point. The path \(X=\zeta+t\) and the square root
are continued from positive \(\zeta\) to the optical branch.
The simple-turning-point multiplier at the root has its universal
Airy value. Expressing the incoming and outgoing modes in the
normalization at infinity multiplies their ratio by the exponential
of minus twice the action shift in \eqref{eq:lovelock-displaced-action}.
Thus the leading WKB term is \(D_{\rm A}(\zeta)\sim\pi/\sqrt\zeta\),
and \eqref{eq:lovelock-crossover} gives
\begin{equation}
 \mathcal B(\chi)\sim\frac{1+\ii}{2\chi}.
 \label{eq:lovelock-massive-limit}
\end{equation}
For example, at \(\chi=5\) the Airy quadrature gives
\(\chi\mathcal B(\chi)=0.50139+0.50137\ii\).
For \(\chi=10\) and \(n=256\), the exact pole slope gives
\(\chi N^{2/3}b_n\simeq0.500019(1+\ii)\).

The same momentum derivative follows from the optical geodesic.
In this metric the inner-to-outer cycle has
\(J_{\rm opt}=\pi\), as in \eqref{eq:jacobi-ads5}.
If \(M\simeq\nu\) is the large worldline mass and \(k=MP\),
then the optical cycle action satisfies
\(\partial_{P^2}W_{\rm opt}=-J_{\rm opt}/2\).
Consequently the logarithmic \(k^2\) derivative of
\(\exp(-MW_{\rm opt})\) equals \(\pi/(2M)\). In the
pole condition with optical phase \(2\omega\mathcal Z_{\rm L}\)
this yields \(b_n\sim\ii\pi/(4M\mathcal Z_{\rm L})=(1+\ii)/(2M)\),
in agreement with \eqref{eq:lovelock-massive-limit}.
For \(|E|=|\omega|/M\), the condition \(M|E|^{-2}\gg1\) for the
geodesic approximation is \(\chi^3/8\gg1\) in this limit.
Each additional optical reflection adds the same quadratic
momentum term in the exponent. Its Gaussian integral has the
transverse multiplicity \(d-1\) and the repetition dependence
of \eqref{eq:jacobi-ratio}. The comparison fixes this part of the
amplitude in the overlap regime. The full correlator also requires
the residues and contour weights. Section~\ref{sec:lovelock-residues}
determines the residue shifts at first order in \(k^2\).
